\chapter{Continuously measured open quantum systems}\noindent
In the previous chapters, we went through the steps leading to the Lindblad master equation using the input-output formalism. The environment has then been traced out to recover information on the system $S$ from the global coupled system, leaving an impact on the dynamics of $S$ crystallized in the expression of a dissipator $\mathfrak{D}$. In the following, we will continuously monitor the environment by measuring the environmental quantum oscillator. Moreover, two kind of measurements will be presented below: namely the photodetection and the homodyne detection. In each case, the reader will be able to find graphical illustrations accompanying the physical equations, which have been generated in \verb|Julia| using the packages \verb|QuantumToolbox.jl| to simulate dynamics, and \verb|Plots.jl| to generate graphs.

We should now calculate the expression of the conditional state obtained at time $t+\dt$. Let $\qty{\ketbra{j}_t}_j$ be a PVM (see Section~\ref{sec:qmt}). The measurement projects the environment with corresponding oscillator $B_t$ on the PVM element $\ketbra{j}_t$. After the measurement on the total state and after applying the partial trace $\mathrm{Tr}_E$, we find
\begin{equation}
    \rho_j(t+\dt)=\frac{\bra{j}_t\bm{\varrho}(t+\dt)\ket{j}_t}{p_j(t+\dt)},
\end{equation}
where $p_j(t+\dt)=\Tr{\tilde{\bm{\varrho}}(t+\dt)\qty(\mathbb{I}\otimes\ketbra{j}_t)}$. Note that in this expression 
\begin{equation}
    \tilde{\bm{\varrho}}(t+\dt)=U_I(t+\dt,t)\overbrace{\rho^{(c)}(t)\otimes\ketbra{0}_t}U_I^\dagger(t+\dt,t),
\end{equation}
where the superscript $(c)$ indicates the system state has been conditioned. If the conditioning is obvious, we choose not to write this superscript explicitly, instead we write for example $\rho_j$.
\section{Photodetection}
\subsection[Deriving the SME]{Deriving the Stochastic Master Equation (SME)}\noindent
In the case where a system is continuously monitored by a photo-detector, the mathematical impact is a projection of the environment on the corresponding Fock basis $\qty{\ketbra{n}_t}_{n\in\mathbb{N}}$. The Eq.~(\ref{eq:expansion}) shows that only projections on $\ketbra{0}_t$ and $\ketbra{1}_t$ contribute to the result. We will show this separately for each measurement result ($0$ or $1$). In order to lighten notations, we will not specify with $\tilde{\bullet}$ that we work in the interaction picture anymore, but if we change the frame of representation it will be specified; also, we work from now on in natural units where $\hbar=1$.
\subsubsection{Measurement result 0:}\noindent
From the expansion in Eq.~(\ref{eq:expansion}), we find
\begin{equation}
    \rho_0(t+\dt)=\frac{\bra{0}_t\bm{\varrho}(t+\dt)\ket{0}_t}{p_0(t+\dt)}= \frac{\rho^{(c)}(t)-\qty{\rho^{(c)}(t),c^\dagger c}\frac{k\dt}{2}}{p_0(t+\dt)}.
\end{equation}
The same expansion allows to calculate 
\begin{equation}
    p_0(t+\dt)=1-\expval{c^\dagger c}_t k\dt,
\end{equation}
where we defined the average labelled by $t$ as $\expval{\bullet}_t=\Tr{\rho^{(c)}(t)\bullet}$. In the infinitesimal limit, we can also expand the denominator in the evolved conditional state, giving
\begin{equation}
    \rho_0(t+\dt)\approx\bar{\rho}_0(t+\dt)\times\qty(1-\expval{c^\dagger c}_t k\dt),
\end{equation}
note that we used the same notation as in Section~\ref{sec:qmt}, where $\bar{\rho}$ indicates an unnormalized state (here in the interaction picture in particular). Finally, by introducing the measurement super-operator $\mathfrak{H}[\hat{O}]\bullet=\hat{O}\bullet + \bullet \hat{O}^\dagger - \expval{\hat{O}+\hat{O}^\dagger}_t\bullet$, we can write that
\begin{equation}\label{eq:rho0}
    \rho_0(t+\dt) = \rho^{(c)}(t)-\frac{k\dt}{2}\mathfrak{H}\qty[c^\dagger c]\rho^{(c)}(t).
\end{equation}
\subsubsection{Measurement result 1:}\noindent
The same calculation method applies and the conditional (unnormalized) state is expressed as
\begin{equation}
    \bar{\rho}_1(t)=\bra{1}_t\bm{\varrho}(t+dt)\ket{1}_t=c\rho^{(c)}(t)c^\dagger k\dt.
\end{equation}
The normalization is then
\begin{equation}
    p_1(t+\dt)=\expval{c^\dagger c}_t k\dt.
\end{equation}
Combined, they give the normalized conditional state
\begin{equation}\label{eq:rho1}
    \rho_1(t+\dt)=\frac{c\rho^{(c)}(t)c^\dagger}{\expval{c^\dagger c}_t}.
\end{equation} 
\subsubsection{Mixing the results:}\noindent
If we observe from a statistical point of view the last results, we have a measurement outcome at time $t+\dt$ that can be $0$ with a probability $p_0(t+\dt)$ or $1$ with probability $p_1(t+\dt)$. We can imagine a process where we count how many times the result $1$ appears, storing the result in a random variable $N(t,t_0)$. In an infinitesimal slice of time $\dt$, the increment $\mathrm{d}N_t=N(t+\dt,t_0)-N(t,t_0)$ is either $1$ or $0$, and so we have the fundamental relation $\mathrm{d}N_t^2=\mathrm{d}N_t$. These are respectively called Poisson process and Poisson increment\footnotemark[1]. For now, we can calculate the stochastic average of $\mathrm{d}N_t$ to get
\begin{equation}\label{eq:avPoisson}
    \mathbb{E}\qty[\mathrm{d}N_t]=0\times p_0(t+\dt)+1\times p_1(t+\dt) = \expval{c^\dagger c}_t k\dt.
\end{equation}
\footnotetext[1]{The Poisson increment follows a Poisson distribution with rate $\Lambda_t\in\mathbb{R}$: $\forall n\in\mathbb{N},\quad\mathbb{P}\qty[\mathrm{d}N_t=n]=\frac{\Lambda_t^n}{n!}e^{-\Lambda_t}$. However because we only have possible outcomes (0 or 1), the parameter is forced to be $\Lambda_t=0$ in our case. We can state that $\mathbb{P}\qty[\mathrm{d}N_t=1]=\lambda_t\dt\approx 0$ and $\mathbb{P}\qty[\mathrm{d}N_t=0]=1-\lambda_t\dt\approx 1$ to obtain the Poisson distribution results on an infinitesimal scale. It is trivial that $\mathbb{E}\qty[\mathrm{d}N_t]=\lambda_t\dt$, which corresponds to results found in Eq~(\ref{eq:avPoisson}).} 
This Poisson increment will allow us to write the differential conditional state $\mathrm{d}\rho^{(c)}(t+\dt)$ with respect to each possible result, because $\mathrm{d}N_t=0$ means the differential is $\mathrm{d}\rho_0(t)=\rho_0(t+\dt)-\rho^{(c)}(t)$, instead $\mathrm{d}N_t=1$ means we the differential is $\mathrm{d}\rho_1(t)=\rho_1(t+\dt)-\rho^{(c)}(t)$ state. Indeed, 
\begin{equation}
    \mathrm{d}\rho^{(c)}_t=\mathrm{d}N_t\qty[\rho_1(t+\dt)-\rho^{(c)}(t)]+\qty(1-\mathrm{d}N_t)\qty[\rho_0(t+\dt)-\rho^{(c)}(t)].
\end{equation}
We can inject in this last equation the results of Eq.~(\ref{eq:rho0},\ref{eq:rho1}) to get (in the Schrödinger picture),
\begin{equation}
    \boxed{\mathrm{d}\rho^{(c)}=-i\qty[H_0,\rho^{(c)}]\dt-\frac{k}{2}\mathfrak{H}\qty[c^\dagger c]\rho^{(c)}\dt + \qty(\frac{c\rho^{(c)}c^\dagger}{\expval{c^\dagger c}_t}-\rho^{(c)})\mathrm{d}N_t},
\end{equation}
in which we discarded terms of order $\mathrm{d}N_t\dt$ (this can be understood informally because $\mathbb{E}\qty[\mathrm{d}N_t]=\lambda_t\dt$ and $\mathbb{V}\qty[\mathrm{d}N_t]\approx\lambda_t\dt$). This equation is called the photo-detection Stochastic Master Equation (SME). The importance of this equation is double because it contains in average the \textbf{unconditional} dynamics: $\rho(t)=\mathbb{E}\qty[\rho^{(c)}(t)]$. Physically, averaging over all possible conditional trajectories gives the unconditional trajectory (as if no measurement was ever done). We can show this by using the average in Eq.~(\ref{eq:avPoisson}) to find
\begin{equation}
    \mathbb{E}\qty[\mathrm{d}\rho^{(c)}]=-i\qty[H_0,\rho]\dt+k\mathfrak{D}[c]\rho\dt,
\end{equation}
which is exactly the Lindblad master equation. An important observation is that doing this entire SME derivation using a pure state instead of a density matrix leads to the following Stochastic Schrödinger Equation (SSE):
\begin{equation}
    \mathrm{d}\ket{\psi^{(c)}}=\qty[-iH_0\dt+\frac{k}{2}\qty(\expval{c^\dagger c}_t-c^\dagger c)\dt + \qty(\frac{c}{\sqrt{\expval{c^\dagger c}_t}}-1)\mathrm{d}N_t]\ket{\psi^{(c)}},
\end{equation}
which conserves the purity of the state during the stochastic evolution. Note that we can derive the SME from the SSE using the Ito calculus relation
\begin{equation}
    \mathrm{d}\rho^{(c)}=\mathrm{d}\ket{\psi^{(c)}}\bra{\psi^{(c)}}
    +\ket{\psi^{(c)}}\mathrm{d}\bra{\psi^{(c)}}
    +\mathrm{d}\ket{\psi^{(c)}}\mathrm{d}\bra{\psi^{(c)}}.
\end{equation}
\section{Homodyne detection}

\section{Linear quantum trajectories}\label{sec:Linear}